\subsubsection{\'{E}tude d'un système ternaire Li-Co-O}
\label{subsubsection:etude_dun_systeme_ternaire_LiCoO}

\setcounter{equation}{0}
\setcounter{Q}{0}

Dans les années 80, les \'{E}tats-Unis ont lancé un large programme de développement de batteries pour la propulsion de véhicules à l'aide de batteries composées d'un alliage de lithium au pôle négatif et de sulfures métalliques au pôle positif. 

Durant ce programme de développement, de nombreuses études thermodynamiques ont été réalisées sur les cathodes dérivées d'oxyde métallique. En particulier, tout un travail a été conduit sur l'étude par titrage coulométrique à \SI{400}{\celsius} des oxydes de cobalt à l'aide de la pile suivante :

\begin{align*}
\ominus ~  Al(s) | \ce{LiAl(s)} | \ce{LiCl}-\ce{KCl} | \ce{Li_xCo_yO_z} ~ \oplus 
\end{align*}

Lors des expériences de titrage coulométrique, la présence des oxydes métalliques suivants : \ce{LiCoO2}, \ce{Co3O4} \ce{CoO} et \ce{Li2O} a été confirmé par diffraction des rayons X.

L'alliage lithium-aluminium \ce{LiAl(s)} a un potentiel connu et mesuré à \SI{400}{\celsius} qui vaut \SI{326}{\milli\volt\per{\ce{Li+ / Li}}}.
%value from paper Gallium, Huggins

\paragraph*{Identification des lignes de démarcation du diagramme ternaire de \ce{LiCoO2} à \SI{400}{\celsius}}

~\

\begin{minipage}{0.48\textwidth}
\includegraphics[width=\textwidth]{Ex29/Fig01}
\end{minipage}\hfill\begin{minipage}{0.48\textwidth}
\captionof{figure}{Diagramme ternaire à compléter.}
\label{fig:29-01}
\end{minipage}

\Q Dans la Fig. \ref{fig:29-01}, placer les corps purs.
\solution{On cherche les coordonnées $(x, y, z)$ à partir de la st{\oe}chiométrie des phases afin de placer les corps purs.
\begin{center}
\begin{tabular}{cccc}
\text{Phase} & $x$ & $y$ & $z$ \\
\ce{O} & 0 & 0 & 1 \\
\ce{Co} & 0 & 1 & 0 \\
\ce{Li} & 1 & 0 & 0 \\
\ce{Li2O} & $\frac{2}{3}$ & 0 & $\frac{1}{3}$ \\
\ce{LiCoO2} & $\frac{1}{4}$ & $\frac{1}{4}$ & $\frac{1}{2}$ \\
\ce{CoO} & $0$ & $\frac{1}{2}$ & $\frac{1}{2}$ \\
\ce{Co3O4} & $0$ & $\frac{3}{7}$ & $\frac{4}{7}$ \\
\end{tabular}
\end{center}
On place alors ces espèces au sein du triangle isotherme de Gibbs : 
\begin{center}
\includegraphics[width=0.48\textwidth]{Ex29/Corr01}
\end{center}
}

Deux titrages coulométriques sur les piles déchargées suivantes :
\begin{align}
\ominus ~  Al(s) | \ce{LiAl(s)} |& \ce{LiCl}-\ce{KCl} | \ce{LiCoO2} ~ \oplus \label{eq:29-pile-1} \\
\ominus ~  Al(s) | \ce{LiAl(s)} |& \ce{LiCl}-\ce{KCl} | \ce{Co3O4} ~ \oplus \label{eq:29-pile-2}
\end{align} 

Les valeurs suivantes de f.e.m. ont été trouvés pour chacun des titrages : 

\begin{figure}[H]
\begin{centering}
\subfloat[\label{fig:29-02a}]{\includegraphics[width=0.48\textwidth]{Ex29/Fig02a}}\hfill
\subfloat[\label{fig:29-02b}]{\includegraphics[width=0.48\textwidth]{Ex29/Fig02b}}
\end{centering}
\caption{\'{E}volution de la force électromotrice par rapport à la référence \ce{Li^+}/\ce{LiAl} à \SI{400}{\celsius} \ref{fig:29-02a} dans le cas d'un pôle $\oplus$ constitué à l'état chargé d'une phase \ce{LiCoO2} ; \ref{fig:29-02b} dans le cas d'un pôle $\oplus$ constitué à l'état chargé d'une phase \ce{CoO}.}
\end{figure}

où $x$ indique le pourcentage molaire de lithium dans la cathode.



\Q Corriger les potentiels pour vous ramener à la demi-pile de référence \ce{Li^+}|\ce{Li(s)} des batteries lithium.
\solution{On va devoir décaler l'échelle des potentiels rédox. Pour cela, on rappelle la définition de la f.e.m. mesurée $e^\circ$ dans notre système :
\begin{align*}
e^\circ &= E_\oplus^\circ - \underbrace{E^\circ_\ominus}_{LiAl/Li} \\
&= \left( E_\oplus^\circ - \underbrace{E'^\circ_\ominus}_{Li^+ / Li} \right) - \left( \underbrace{E^\circ_\ominus}_{Li^+/LiAl} - \underbrace{E'^\circ_\ominus}_{Li^+ / Li} \right) \\
e^\circ &= e'^\circ - 0.326
\end{align*}
où $e'^\circ$ serait le potentiel mesuré dans une pile où le pôle $\ominus$ est composé de lithium métallique. \\
\textbf{AN} :
\begin{center}
\begin{tabular}{cc}
\toprule
$e^\circ$ & $e'^\circ$ \\
\si{\volt} & \si{\volt} \\
\midrule
1.310 & 1.636 \\
-0.326 & 0.0 \\
1.481 & 1.807 \\
\bottomrule
\end{tabular}
\end{center}}

\Q Que signifie le plateau ? Combien de transition de phase observe-t-on dans chacune des piles ?
\solution{
Un plateau signifie que tous les paramètres intensifs du système sont imposés, c'est-à-dire que la variance particulière $v_p = 0$. Dans un système ternaire (c'est-à-dire qu'il y a 3 constituants chimiques soit $N=3$), sachant que la pression et la température sont imposées dans ces expériences, on peut alors écrire : 
\begin{align*}
v_p &= v - 2 \\
&= (2 + N - P) - 2 = N - P \\
&\implies P = N - v_p \stackrel{AN}{=} 0 
\end{align*}
Un plateau traduit la présence d'un mélange hétérogène constitué de 3 phases non miscibles.}

\Q En représentant le segment où le lithium varie de composition dans une pile pour les deux pôles $\oplus$, indiquer sur ces segments les points de transition de phase. En déduire les domaines ternaires du mélange.
\solution{
On représente les segments entre les pôles purs {\color{blue} \ce{Li}-\ce{LiCoO2}} et {\color{red} \ce{Li}-\ce{CoO}}.
\begin{center}
\begin{minipage}{0.48\textwidth}
\includegraphics[width=\textwidth]{Ex29/Corr02}
\end{minipage}\hfill\begin{minipage}{0.48\textwidth}
\includegraphics[width=\textwidth]{Ex29/Corr03}
\end{minipage}
\end{center}
On peut alors facilement intuiter en s'aidant des valeurs de f.e.m. Les valeurs à \SI{0.00}{\volt\per{Li^+ / Li}} et \SI{1.636}{\volt\per{Li^+ / Li}} sont communes aux deux droites et indiquent que sur chacun de ces segments ont les trois mêmes phases dans le mélange. On peut alors directement intuiter la ligne de démarcation entre des deux domaines. L'autre ligne de démarcation doit nécessairement passer par le pôle \ce{LiCoO2} et par la transition \SI{1.807}{\volt\per{Li^+ / Li}} - \SI{1.636}{\volt\per{Li^+ / Li}}.}

\paragraph*{Relation entre stabilité des oxydes non lithiés et potentiel d'électrode}

Lors de cette étude à \SI{400}{\celsius}, les chercheurs ont élargi l'étude précédente à l'étude d'autres ternaires : \ce{Li}-\ce{Mn}-\ce{O}, \ce{Li}-\ce{Fe}-\ce{O} et \ce{Li}-\ce{Co}-\ce{O}. 

On va chercher à étudier un sous-triangle composé de deux oxydes : \ce{MOy}, \ce{MOy-x} et de l'oxyde de lithium \ce{Li2O} et on considérera que ces oxydes ne sont pas lithiés. Lors du développement du modèle, les corps purs auront une activité de 1 puisqu'ils sont seuls dans leur phase et le dioxygène sera assimilé à un gaz parfait.

Dans un sous triangle de ce type, on étudie l'équilibre général suivant : 
\begin{equation}
\ce{MOy + 2x Li = MO_{y-x} + x Li2O} \label{eq:29-01}
\end{equation}

où l'enthalpie libre de réaction est donnée par : 
\begin{align*}
\Delta_r G^\circ = -2 x F e^\circ
\end{align*}

où $e^\circ$ désigne maintenant la force électromotrice par rapport à une référence lithium métallique.

\Q Identifier un sous-triangle du diagramme \ce{Li}-\ce{Co}-\ce{O} de la Fig. \ref{fig:29-01} qui pourra être comparé à un système équivalent non lithié.
\solution{Le sous-triangle composé de \ce{Li2O}-\ce{LiCoO2}-\ce{Co} a été étudié sous la forme \ce{Li2O}-\ce{CoO2}-\ce{Co}.}

\Q \'{E}crire l'équation de corrosion sèche des couples \ce{MO_{y-x}}/\ce{MO_y} \& \ce{Li2O}/\ce{Li}. En déduire une équation de réaction.
\solution{
On pose deux équations d'oxydation sèche par \ce{O2}.
\begin{center}
\begin{minipage}{0.4\textwidth}
\begin{align*}
\ce{MOy &= \frac{x}{2} O2 + MO_{y-x}} \quad \Delta_r G^\circ ( \ce{MOy} ) \\ 
\ce{2x Li + \frac{x}{2} O2 &= x Li2O}  \quad \Delta_r G^\circ (\ce{Li2O}) \\
\hline
\ce{MOy + 2x Li &= MO_{y-x} + x Li2O} \quad \Delta_r G^\circ
\end{align*}
\end{minipage}
\end{center}}

\Q Exprimer l'enthalpie standard de réaction de corrosion sèche de  \ce{Li2O}/\ce{Li} en fonction de son enthalpie de formation, notée $\Delta_f G^\circ ( \ce{Li2O})$.
\solution{
La réaction de formation de \ce{Li2O} est la réaction suivante : 
\begin{align*}
\ce{2Li (s) + \frac{1}{2} O2 (g)= Li2O(s)}
\end{align*}
On a donc immédiatement que : 
\begin{align*}
\Delta_r G^\circ (\ce{Li2O}) = x \Delta_f G^\circ ( \ce{Li2O} )
\end{align*}}

\Q Exprimer l'enthalpie standard de réaction de corrosion sèche de \ce{MO_{y-x}}/\ce{MO_y} en fonction de la pression partielle d'équilibre en \ce{O2}, notée $P_{eq}(\ce{O2})$.
\solution{On sait que : 
\begin{align*}
\Delta_r G^\circ ( \ce{MOy} ) = -RT \ln K_r^\circ
\end{align*}
Avec la loi d'action des masses, on sait que : 
\begin{align*}
K^\circ (T) &\stackrel{eq}{=} \frac{a (\ce{MO_{y-x}}) a^\frac{x}{2}(\ce{O2})}{a (\ce{MO_{y}})} \\
&= \left( \frac{P_{eq} ( \ce{O2} )}{P^\circ} \right)^\frac{x}{2}
\end{align*}
En combinant les deux, on a alors : 
\begin{align*}
\Delta_r G^\circ ( \ce{MOy} ) = -\frac{xRT}{2} \ln \frac{P_{eq} ( \ce{O2} )}{P^\circ}
\end{align*}}

\Q En déduire la relation entre la force électromotrice et les paramètres : $\Delta_f G^\circ ( \ce{Li2O})$ et $P_{eq}(\ce{O2})$.
\solution{\`{A} partir de toutes les relations précédentes, on sait que : 
\begin{align*}
\Delta_r G^\circ &=  \Delta_r G^\circ ( \ce{MOy} ) + \Delta_r G^\circ ( \ce{Li2O} ) \\
-2x F e^\circ &=  x \Delta_f G ( \ce{MOy} ) - \frac{xRT}{2} \ln \frac{P_{eq} ( \ce{O2} )}{P^\circ} \\
e^\circ &= \frac{RT}{4F} \ln P_{eq} ( \ce{O2} ) - \frac{\Delta_f G^\circ ( \ce{Li2O} ) }{2F}
\end{align*}
}

Dans divers systèmes lithiés, la force électromotrice est mesurée et mise en relation avec la pression partielle des diagrammes d'Ellingham (cf cours). Il est observé une quasi-parfaite corrélation entre les résultats des diagrammes et la pression en \ce{O2} : 
\begin{center}
\begin{minipage}{0.48\textwidth}
\includegraphics[width=\textwidth]{Ex29/Fig03}
\end{minipage}\hfill\begin{minipage}{0.48\textwidth}
\captionof{figure}{\'{E}volution du potentiel du pôle $\oplus$ composés de systèmes ternaires lithiés en fonction de la pression d'équilibre tirée des diagrammes d'Ellingham des oxydes non lithiés.}
\end{minipage}
\end{center}

L'équation de la droite de corrélation à \SI{400}{\celsius} est : 
\begin{align}
E_\oplus^\circ = 3.34 \cdot 10^{-2} \log P ( \ce{O2} ) + 2.65  \quad [\si{\volt\per{\ce{Li^+} / \ce{Li}}}] 
\label{eq:29-02}
\end{align}

Il a été montré par la suite que cette relation est vérifiée aussi à d'autres températures (\emph{par ex.} les conditions ambiantes) et pour tout système contenant deux oxydes binaire et l'oxyde de lithium.

%De plus, on sait que les matériaux se comporte différemment en fonction de leur potentiel rédox. \`{A} grand potentiel, les matériaux s'écartent de la st{\oe}chiométrie et présentent des lacunes électroniques. \`{A} faible potentiel, les matériaux présentent plutôt un excès d'électron.

\Q Commenter si la corrélation positive a du sens. Quel intérêt présente une telle corrélation ? 
\solution{La corrélation nous montre qu'il est plus facile de lithier les matériaux qui présentent la plus inertie chimique à l'oxydation sèche. Au bilan, les matériaux les plus "inerte" chimique vis-à-vis du \ce{O2} sont les plus capables d'accepter du lithium. \\
Cette relation suppose que la connaissance du comportement thermique des oxydes non lithié permet de prédire le potentiel de la cathode ! Un autre avantage de cette relation est que cette relation est indépendante de considérations structurales sur l'oxyde en question. \\
Bien entendu, il existe des déviations à cette relation en particulier quand la stabilité de l'oxyde est sensible à la présence du lithium ! 
}

\Q La pression partielle d'équilibre en \ce{O2} du sous-triangle \ce{Li2O}-\ce{LiCoO2}-\ce{Co} est attendue autour de $10^{-30}$ \si{\bar}. Calculer le potentiel attendu et comparer là à la valeur obtenue expérimentalement de \SI{1.807}{\volt\per{\ce{Li^+}/\ce{Li}}}.
\solution{On obtient une tension de \SI{1.684}{\volt\per{\ce{Li^+}/\ce{Li}}} à l'aide de l'Eq. \ref{eq:29-02}. \\
Or ici, on compare deux ternaires différents :
\begin{center}
\begin{tabular}{cc}
\ce{Li2O}-\ce{LiCoO2}-\ce{Co} & \SI{1.807}{\volt\per{\ce{Li^+}/\ce{Li}}} \\
\ce{Li2O}-\ce{CoO2}-\ce{Co} & \SI{1.684}{\volt\per{\ce{Li^+}/\ce{Li}}}
\end{tabular}
\end{center}
Cela traduit le fait que \ce{LiCoO2} est plus stable thermodynamiquement que le composé \ce{CoO2} (c-à-d que le dérivé du cobalt préfère être lithié) ! Un beau résultat de thermodynamique à mettre en perspective avec nos études cristallographiques à la section suivante.}


